\section{Apparatus and Experimental Techniques}
\label{sec:apparatus}

% PAGE TARGET: ~4 pages.
% Day 3 (5/27) afternoon -- right after Theory.
% CUED: describe equipment and runs but DO NOT give a blow-by-blow.
% Discuss experimental accuracy here.

\subsection{Imaging Equipment}
\label{sec:apparatus:hw}
% - MSI rig used at \emph{TODO: library name}: bands, sensor, FOV.
% - DRP acquisition: lighting dome, $N_\theta \times N_\phi$ sampling,
%   exposure per frame.
% - TX (transmitted-light) acquisition: light box specification.

\subsection{Manuscript Dataset}
\label{sec:apparatus:dataset}
% - 9 folios across 4 manuscript stocks: Kk1-5, Hh2-10/-12, Ee5-22,
%   Ff2-6 / Ff4-9 / Ff4-15, Ii3-8 (Cambridge UL shelfmarks).
% - Per-folio acquisition table (FOV cm, ROI crop, modality).
% - Spreadsheet GT source + caveats (stock vs folio).
% - Manual GT process for Kk1-5 f5v / f9v.

\subsection{Synthetic Phantom Generator}
\label{sec:apparatus:phantom}
% Module: scripts/phantom\_synthetic.py.
% - Periodic Gaussian-profile lines under controllable
%   (period, sigma, angle, contrast, SNR).
% - 5 random seeds per sweep point.
% - This is the testbed for Section 4.1 sweeps.

\subsection{Software Implementation}
\label{sec:apparatus:sw}
% - Python 3.10+; numpy, scipy, scikit-image, OpenCV.
% - 5-stage modular layout under paperdrm/.
% - DRP stack cached as numpy memmap; cache invalidation by serial+shape.
% - Run on commodity laptop hardware (specify in Appendix B).

\subsection{Calibration and Experimental Accuracy}
\label{sec:apparatus:cal}
% - FOV scale set by physical ruler in image; uncertainty $\sim$0.5\%.
% - Background subtraction reference: Gaussian-blurred mean of stack.
% - Repeat measurement spread: see split-half stability (Section 4.2).
% - Polarity (wire\_is\_darker) set per modality; documented per-folio
%   exp\_param.yaml.
